\documentclass[aspectratio=169]{beamer}

% --- Step Function Branding (replaces any generic Beamer theme) ---
% DO NOT use \usetheme{...}; sf-branding handles all theming.
\usepackage{sf-branding}

% --- Additional Packages ---
\usepackage{appendixnumberbeamer}
\usepackage{amsmath,amssymb}
\usepackage{siunitx}
\sisetup{group-separator={,}, group-minimum-digits=4}
\usepackage{graphicx}
\usepackage{booktabs}

% --- Project-specific colors (data viz only) ---
\definecolor{revenue}{HTML}{2E7D32}
\definecolor{cost}{HTML}{C62828}

\title{Techno-Economic Assessment of\\[4pt] \textbf{Steel Slag} Valorization}
\subtitle{H\textsubscript{2} Production $\cdot$ CO\textsubscript{2} Sequestration $\cdot$ Iron Recovery}
\author{Step Function LLC}
\date{\today}

\begin{document}

% ═══════════════════════════════════════════════════════════════
\begin{frame}
\titlepage
\end{frame}

% ═══════════════════════════════════════════════════════════════
\begin{frame}{Outline}
\tableofcontents
\end{frame}

% ═══════════════════════════════════════════════════════════════
\section{Background \& Opportunity}
% ═══════════════════════════════════════════════════════════════

\begin{frame}{Opportunity: BOF Steel Slag}
\framesub{High-volume industrial byproduct with untapped value}
\begin{columns}
\column{0.55\textwidth}
\begin{itemize}
    \item \textbf{250+ Mt/yr} slag produced globally
    \item Mostly landfilled — environmental liability
    \item Rich in CaO, FeO/Fe\textsubscript{2}O\textsubscript{3}, MgO, SiO\textsubscript{2}
\end{itemize}
\vspace{8pt}
\textbf{Three valorization pathways:}
\begin{enumerate}
    \item Green H\textsubscript{2} (serpentinization)
    \item CO\textsubscript{2} sequestration (carbonation)
    \item Iron concentrate (magnetic sep.)
\end{enumerate}

\column{0.42\textwidth}
\begin{table}
\centering
\small
\begin{tabular}{@{}lS[table-format=2.0]@{}}
\toprule
\textbf{\textcolor{white}{Component}} & \textbf{\textcolor{white}{wt\%}} \\
\midrule
\textcolor{white}{CaO}                & 40 \\
\textcolor{white}{Fe\textsubscript{2}O\textsubscript{3}} & 15 \\
\textcolor{white}{FeO}                & 15 \\
\textcolor{white}{SiO\textsubscript{2}} & 20 \\
\textcolor{white}{MgO}                & 10 \\
\bottomrule
\end{tabular}
\end{table}
\end{columns}
\end{frame}

% ═══════════════════════════════════════════════════════════════
\begin{frame}{Resource Value Proposition}
\framesub{A circular-economy approach to industrial waste}
\begin{itemize}
    \item \textbf{Environmental}: Avoids landfill; sequesters CO\textsubscript{2} in stable mineral carbonates
    \item \textbf{Economic}: Three revenue streams from a zero-cost feedstock
    \item \textbf{Strategic}: Green H\textsubscript{2} production at industrial scale
\end{itemize}
\vspace{12pt}
\centering
\begin{tabular}{@{}lll@{}}
\toprule
\textbf{\textcolor{white}{Pathway}} & \textbf{\textcolor{white}{Product}} & \textbf{\textcolor{white}{Market Value}} \\
\midrule
Magnetic separation & Fe concentrate & \SI{120}{\$/t} \\
Mineral carbonation & CaCO\textsubscript{3} + CO\textsubscript{2} credits & \SI{60}{\$/t_{CO_2}} \\
Serpentinization    & Green H\textsubscript{2} & \SI{4000}{\$/t} \\
\bottomrule
\end{tabular}
\end{frame}

% ═══════════════════════════════════════════════════════════════
\section{Process Design}
% ═══════════════════════════════════════════════════════════════

\begin{frame}{Process Flow Diagram}
\framesub{Optimal topology: heat recovery only (no aux heater, no scrubber)}
\begin{center}
\begin{tikzpicture}[
  scale=0.82, every node/.append style={transform shape},
  node distance=0.8cm and 1.3cm,
  unit/.style={rectangle, draw=StepGrey, rounded corners=3pt, minimum width=1.6cm,
               minimum height=0.60cm, align=center, fill=StepElectric!15,
               font=\scriptsize, text=white},
  product/.style={rectangle, draw=StepGrey, rounded corners=3pt, minimum width=1.4cm,
                  minimum height=0.50cm, align=center, fill=StepGreen!15,
                  font=\scriptsize\itshape, text=white},
  feed/.style={rectangle, draw=StepGrey, rounded corners=3pt, minimum width=1.4cm,
               minimum height=0.50cm, align=center, fill=StepOrange!15,
               font=\scriptsize\itshape, text=white},
  stream/.style={-{Stealth[length=2mm]}, semithick, white},
]

\node[feed] (slag) {Slag\\(\SI{500}{t/d})};
\node[unit, right=of slag] (lims) {LIMS};
\node[product, above=0.5cm of lims] (fe) {Fe Conc.};
\node[unit, right=of lims] (pump1) {P-101};
\node[unit, right=of pump1] (hx1) {E-101};
\node[unit, right=of hx1] (serp) {R-101\\(\SI{250}{\degreeCelsius})};
\node[unit, below=0.7cm of serp] (sep) {S-101};
\node[product, right=of sep] (h2) {H\textsubscript{2}};

\node[feed, below=1.5cm of pump1] (co2) {CO\textsubscript{2}};
\node[unit, right=of co2] (hx2) {E-102};
\node[unit, right=of hx2] (carb) {R-201\\(\SI{185}{\degreeCelsius})};
\node[unit, below=0.7cm of carb] (thick) {S-201};
\node[product, right=of thick] (caco3) {CaCO\textsubscript{3}};

\draw[stream] (slag) -- (lims);
\draw[stream] (lims) -- (fe);
\draw[stream] (lims) -- (pump1);
\draw[stream] (pump1) -- (hx1);
\draw[stream] (hx1) -- (serp);
\draw[stream] (serp) -- (sep);
\draw[stream] (sep) -- (h2);
\draw[stream] (sep.south) -- ++(0,-0.3) -| (hx2);
\draw[stream] (co2) -- (hx2);
\draw[stream] (hx2) -- (carb);
\draw[stream] (carb) -- (thick);
\draw[stream] (thick) -- (caco3);

\draw[-{Stealth}, dashed, StepTeal!60] (serp.south west) to[bend right=25]
    node[left, font=\tiny, StepTeal]{heat} (hx1.south);
\draw[-{Stealth}, dashed, StepTeal!60] (carb.south west) to[bend right=25]
    node[left, font=\tiny, StepTeal]{heat} (hx2.south);

\end{tikzpicture}
\end{center}
\vspace{4pt}
\small
\textbf{Key insight}: Iron removal before carbonation raises pH for carbonate stability (Reaktoro-confirmed).
\end{frame}

% ═══════════════════════════════════════════════════════════════
\begin{frame}{Key Reactions}
\framesub{Two complementary conversion pathways}
\begin{block}{Serpentinization (H\textsubscript{2} Production)}
\vspace{-6pt}
\begin{equation*}
    \text{Fe}_2\text{SiO}_4 + 3\,\text{H}_2\text{O}
    \xrightarrow{250\,^\circ\text{C},\;100\,\text{bar}}
    2\,\text{Fe(OH)}_2 + \text{SiO}_2 + \text{H}_2
\end{equation*}
\end{block}
\vspace{4pt}
\begin{block}{Mineral Carbonation (CO\textsubscript{2} Sequestration)}
\vspace{-6pt}
\begin{align*}
    \text{CaO} + \text{CO}_2 &\xrightarrow{\SI{185}{\degreeCelsius},\;\SI{50}{bar}} \text{CaCO}_3 \\
    \text{MgO} + \text{CO}_2 &\longrightarrow \text{MgCO}_3
\end{align*}
\end{block}
\end{frame}

% ═══════════════════════════════════════════════════════════════
\begin{frame}{Thermodynamic Speciation}
\framesub{Reaktoro equilibrium confirms iron removal is essential before carbonation}
\begin{table}
\centering
\small
\begin{tabular}{@{}lccp{5cm}@{}}
\toprule
\textbf{\textcolor{white}{Scenario}} & \textbf{\textcolor{white}{pH}} & \textbf{\textcolor{white}{$E_h$}} & \textbf{\textcolor{white}{Key Finding}} \\
\midrule
\textcolor{white}{Ambient (\SI{25}{\degreeCelsius})} & \textcolor{white}{2.43} & \textcolor{white}{0.656} & \textcolor{StepGrey}{Fe$^{3+}$ hydrolysis; highly acidic} \\
\textcolor{white}{Serpent. (\SI{250}{\degreeCelsius})} & \textcolor{white}{5.60} & \textcolor{white}{0.661} & \textcolor{StepGrey}{H$_2$(aq) produced; 0.7\% eq.\ conv.} \\
\textcolor{white}{Carbonation (+CO$_2$)} & \textcolor{white}{3.20} & \textcolor{white}{0.106} & \textcolor{StepGrey}{$<$0.1\% CO$_2$ conv.\ at low pH} \\
\bottomrule
\end{tabular}
\end{table}
\vspace{10pt}
\centering
\textcolor{StepYellow}{\textbf{Critical}}: Dissolved Fe$^{3+}$ suppresses pH to $\sim$2--3, preventing carbonate precipitation.\\
Iron removal must precede carbonation.
\end{frame}

% ═══════════════════════════════════════════════════════════════
\begin{frame}{Modeling Methodology}
\framesub{Three-tier computational hierarchy}
\begin{columns}
\column{0.55\textwidth}
\begin{enumerate}
    \item \textbf{Reaktoro} speciation
    \begin{itemize}
        \item Gibbs energy minimization
        \item Confirms phase stability \& pH
    \end{itemize}
    \item \textbf{IDAES/GDP} superstructure
    \begin{itemize}
        \item Process synthesis \& screening
        \item Topology enumeration
    \end{itemize}
    \item \textbf{JAX} techno-economics
    \begin{itemize}
        \item Differentiable cost model
        \item Autodiff sensitivity analysis
    \end{itemize}
\end{enumerate}

\column{0.42\textwidth}
\centering
\begin{tikzpicture}[
  box/.style={rectangle, draw=StepGrey, rounded corners=3pt, minimum width=2.8cm,
              minimum height=0.6cm, align=center, font=\scriptsize, text=white},
  arr/.style={-{Stealth[length=2mm]}, semithick, white},
]
\node[box, fill=StepElectric!15] (rkt) {Reaktoro\\Speciation};
\node[box, fill=StepTeal!15, below=0.6cm of rkt] (idaes) {IDAES/GDP\\Process Synthesis};
\node[box, fill=StepGreen!15, below=0.6cm of idaes] (jax) {JAX TEA\\Cost \& Sensitivity};
\draw[arr] (rkt) -- (idaes);
\draw[arr] (idaes) -- (jax);
\end{tikzpicture}
\end{columns}
\end{frame}

% ═══════════════════════════════════════════════════════════════
\section{GDP Analysis}
% ═══════════════════════════════════════════════════════════════

\begin{frame}{Superstructure Overview}
\framesub{Generalized Disjunctive Programming for topology screening}
\begin{columns}
\column{0.55\textwidth}
\begin{itemize}
    \item \textbf{15 unit operations} (12 fixed, 3 optional)
    \item \textbf{1 disjunction}: Aux heater \texttt{XOR} Heat recovery only
    \item \textbf{1 optional}: Inter-reactor scrubber
    \item \textbf{4 feasible topologies}
    \item \textbf{12 continuous} design variables for BoTorch
\end{itemize}

\column{0.42\textwidth}
\small
\begin{table}
\centering
\begin{tabular}{@{}cc@{}}
\toprule
\textbf{\textcolor{white}{Heater}} & \textbf{\textcolor{white}{Scrubber}} \\
\midrule
\textcolor{StepGreen}{\checkmark} & \textcolor{StepGreen}{\checkmark} \\
\textcolor{StepGreen}{\checkmark} & \textcolor{StepGrey}{---} \\
\textcolor{StepGrey}{---} & \textcolor{StepGreen}{\checkmark} \\
\textcolor{StepGrey}{---} & \textcolor{StepGrey}{---} \\
\bottomrule
\end{tabular}
\end{table}
\centering
\textcolor{StepGrey}{\small 4 topologies (2$\times$2)}
\end{columns}
\end{frame}

% ═══════════════════════════════════════════════════════════════
\begin{frame}{Topology Ranking}
\framesub{Heat recovery only is the clear winner}
\begin{table}
\centering
\small
\begin{tabular}{@{}clccccc@{}}
\toprule
\textbf{\textcolor{white}{Rank}} & \textbf{\textcolor{white}{Topology}} & \textbf{\textcolor{white}{CAPEX}}
  & \textbf{\textcolor{white}{OPEX}} & \textbf{\textcolor{white}{EAC}}
  & \textbf{\textcolor{white}{Revenue}} & \textbf{\textcolor{white}{Net}} \\
& & \textcolor{StepGrey}{(\$M)} & \textcolor{StepGrey}{(\$M/yr)} & \textcolor{StepGrey}{(\$M/yr)}
  & \textcolor{StepGrey}{(\$M/yr)} & \textcolor{StepGrey}{(\$M/yr)} \\
\midrule
\textbf{\textcolor{StepTeal}{1}} & \textbf{\textcolor{StepTeal}{HR only}} & \textbf{7.76} & \textbf{2.49} & \textbf{3.28}
  & \textbf{9.53} & \textcolor{StepGreen}{\textbf{6.25}} \\
2 & HR + scrubber & 7.96 & 2.49 & 3.30 & 9.53 & 6.23 \\
3 & Aux heater    & 8.06 & 3.07 & 3.89 & 9.53 & 5.64 \\
4 & Aux + scrub   & 8.26 & 3.07 & 3.91 & 9.53 & 5.62 \\
\bottomrule
\end{tabular}
\end{table}
\vspace{6pt}
\centering
\small Aux heater adds \SI{0.58}{M\$/yr} OPEX with no measurable benefit.
\end{frame}

% ═══════════════════════════════════════════════════════════════
\begin{frame}{Optimal Topology}
\framesub{Heat recovery eliminates auxiliary heater capital and energy cost}
\begin{itemize}
    \item \textbf{Selected}: Heat recovery only — no aux heater, no scrubber
    \item \textbf{CAPEX}: \$\num{7760000} (lowest)
    \item \textbf{EAC}: \$\num{3280000}/yr (lowest)
    \item \textbf{Net annual value}: \textcolor{StepGreen}{\textbf{\$\num{6250000}/yr}} (highest)
\end{itemize}
\vspace{12pt}
\textbf{Key differentiators:}
\begin{itemize}
    \item Exothermic serpentinization/carbonation reactions supply heat to exchangers
    \item Scrubber adds \$\num{200000} CAPEX with negligible impact on purity
    \item Simpler flowsheet reduces maintenance and downtime risk
\end{itemize}
\end{frame}

% ═══════════════════════════════════════════════════════════════
\section{Results}
% ═══════════════════════════════════════════════════════════════

\begin{frame}{Simulation KPIs}
\framesub{Recoveries and conversions from IDAES simulation}
\begin{table}
\centering
\begin{tabular}{@{}lr@{}}
\toprule
\textbf{\textcolor{white}{Metric}} & \textbf{\textcolor{white}{Value}} \\
\midrule
Fe recovery (LIMS) & \SI{85}{\percent} \\
H\textsubscript{2} production & \SI{241}{t/yr} \\
CO\textsubscript{2} sequestered & \num{58575}~t/yr \\
CaCO\textsubscript{3} produced & \num{104598}~t/yr \\
Feedstock throughput & \SI{500}{t/d} $\times$ \SI{330}{d/yr} \\
\bottomrule
\end{tabular}
\end{table}
\end{frame}

% ═══════════════════════════════════════════════════════════════
\begin{frame}{Net Value}
\centering
\vspace{2cm}
{\Huge\textbf{\textcolor{StepGreen}{\$37.87/t}}}\\[12pt]
{\large Net value per tonne of steel slag}\\[8pt]
{\normalsize
\textcolor{StepGreen}{Revenue: \$57.73/t} \quad
\textcolor{StepRed}{Cost: \$19.88/t}}\\[6pt]
{\small\textcolor{StepGrey}{\SI{500}{t/d} $\cdot$ \SI{330}{d/yr} $\cdot$ 20-yr project $\cdot$ 8\% discount}}
\end{frame}

% ═══════════════════════════════════════════════════════════════
\begin{frame}{Revenue Breakdown}
\framesub{Iron concentrate is the dominant value driver}
\begin{table}
\centering
\begin{tabular}{@{}lrrc@{}}
\toprule
\textbf{\textcolor{white}{Product}} & \textbf{\textcolor{white}{Production}} & \textbf{\textcolor{white}{Revenue}} & \textbf{\textcolor{white}{Share}} \\
\midrule
Fe concentrate & \num{42075}~t/yr & \textcolor{StepGreen}{\$5.05M/yr} & 53\% \\
CO\textsubscript{2} credits & \num{58575}~t/yr & \textcolor{StepGreen}{\$3.51M/yr} & 37\% \\
H\textsubscript{2} gas & \SI{241}{t/yr} & \textcolor{StepGreen}{\$0.96M/yr} & 10\% \\
\midrule
\textbf{Total} & & \textcolor{StepGreen}{\textbf{\$9.53M/yr}} & 100\% \\
\bottomrule
\end{tabular}
\end{table}
\vspace{10pt}
\centering
\small Iron concentrate is the dominant value driver at 53\% of total revenue.
\end{frame}

% ═══════════════════════════════════════════════════════════════
\begin{frame}{Economics Summary}
\framesub{Annual cost vs.\ revenue breakdown}
\begin{table}
\centering
\begin{tabular}{@{}lr@{}}
\toprule
\textbf{\textcolor{white}{Item}} & \textbf{\textcolor{white}{Value}} \\
\midrule
\textcolor{StepRed}{CAPEX (installed)} & \$\num{7760000} \\
\textcolor{StepRed}{OPEX (annual)} & \$\num{2490000}/yr \\
\textcolor{StepRed}{EAC (annualized total)} & \$\num{3280000}/yr \\
\midrule
\textcolor{StepGreen}{Revenue (annual)} & \$\num{9530000}/yr \\
\textcolor{StepGreen}{\textbf{Net annual value}} & \textbf{\$\num{6250000}/yr} \\
\bottomrule
\end{tabular}
\end{table}
\vspace{8pt}
\centering
\small Payback period: $<$ 2~years at full capacity.
\end{frame}

% ═══════════════════════════════════════════════════════════════
\section{Optimization \& Sensitivity}
% ═══════════════════════════════════════════════════════════════

\begin{frame}{BoTorch Optimization Results}
\framesub{Bayesian optimization over 12 continuous parameters}
\begin{table}
\centering
\begin{tabular}{@{}lcc@{}}
\toprule
\textbf{\textcolor{white}{Metric}} & \textbf{\textcolor{white}{Baseline}} & \textbf{\textcolor{white}{Optimized}} \\
\midrule
OPEX (\$/yr) & \num{2490000} & \num{2350000} \\
EAC (\$/yr)  & \num{3280000} & \num{3140000} \\
Net (\$/yr)  & \num{6250000} & \textcolor{StepGreen}{\textbf{\num{6390000}}} \\
\midrule
\textbf{Improvement} & & \textcolor{StepGreen}{+\SI{2.2}{\percent}} \\
\bottomrule
\end{tabular}
\end{table}
\vspace{8pt}
\centering
\small 15 BO iterations converged; optimal point reduces energy consumption via lower residence time.
\end{frame}

% ═══════════════════════════════════════════════════════════════
\begin{frame}{Cost Sensitivity (JAX Autodiff)}
\framesub{Top cost drivers identified via automatic differentiation}
\begin{table}
\centering
\begin{tabular}{@{}lrl@{}}
\toprule
\textbf{\textcolor{white}{Parameter}} & \textbf{\textcolor{white}{$|\partial\text{EAC}/\partial x|$}} & \textbf{\textcolor{white}{Interpretation}} \\
\midrule
Residence time & \SI{10185}{\$/h}  & Reactor sizing driver \\
Operating temp & \SI{8250}{\$/\degreeCelsius} & Energy cost driver \\
Throughput     & \SI{5214}{\$/(t/d)} & Economies of scale \\
\bottomrule
\end{tabular}
\end{table}
\vspace{10pt}
\textbf{Key risks:}
\begin{itemize}
    \item CO\textsubscript{2} credits at \SI{30}{\$/t} $\to$ net value drops to \$27.22/t
    \item Fe price at \SI{80}{\$/t} $\to$ revenue loss of \$1.68M/yr
    \item Serpentinization kinetics require catalyst R\&D
\end{itemize}
\end{frame}

% ═══════════════════════════════════════════════════════════════
\section{Conclusions}
% ═══════════════════════════════════════════════════════════════

\begin{frame}{Key Findings}
\begin{enumerate}
    \item Steel slag valorization at \SI{500}{t/d} yields \textcolor{StepGreen}{\textbf{\$37.87/t}} net value
    \item Optimal topology: \textbf{heat recovery only} (no auxiliary heater, no scrubber)
    \item \textbf{Iron concentrate} is the dominant revenue driver (53\%)
    \item \textbf{Iron removal before carbonation} is thermodynamically essential
    \item Aux heater \& scrubber do not justify added cost
\end{enumerate}
\vspace{12pt}
\textbf{Next steps:}
\begin{itemize}
    \item Pilot-scale serpentinization kinetics testing
    \item Detailed vendor quotes for high-pressure reactor vessels
    \item Include Cr/Mn/V value streams (requires expanded thermodynamic database)
\end{itemize}
\end{frame}

% ═══════════════════════════════════════════════════════════════
% Branded closing slide
\closingSlide

\end{document}
